\documentclass[aps,pra,onecolumn,superscriptaddress,nofootinbib]{revtex4-2}
\usepackage{amsmath,amssymb,bm}
\usepackage{booktabs}
\usepackage{siunitx}
\usepackage{hyperref}
\hypersetup{colorlinks=true,citecolor=blue,urlcolor=blue,linkcolor=blue}
\sisetup{detect-all}

\begin{document}

\title{Supplemental Material for: Apparent diffusion from deterministic velocity gradients in wavelength-resolved photodetectors}
\author{Anonymous}
\date{August 16, 2026}
\maketitle

\section{Scope and separation of the two conditional models}

The manuscript uses two distinct numerical constructions that must not be conflated.  First, a one-dimensional monotonic graded-HgCdTe optical model is used only to generate six finite, wavelength-dependent theoretical generation kernels $g_m(z)$.  Those same generated kernel shapes are supplied exactly to the forward averaging and to the kernel-aware inverse.  They are not measured or experimentally calibrated kernels from a specific detector.  Second, the central apparent-diffusion counterexample is generated by a separate deterministic planar-depletion transport stress with microscopic diffusion fixed to zero.  The graded optical-kernel model therefore determines how the point-source response is sampled; it is not the electrostatic transport model that generates the counterexample.

All numerical quantities below are conditional theoretical-model outputs.  They do not constitute a calibrated simulation of a published detector.

\section{Conditional HgCdTe optical-kernel construction}

The optical coordinate is $0\le z\le L$ with $L=7.6\,\mu\mathrm{m}$.  At $T=300\,\mathrm{K}$ the Cd mole fraction is taken to vary monotonically from $x=0.55$ at $z=0$ to $x=0.32$ at $z=L$.  The band gap is evaluated with the Hansen--Schmit--Casselman empirical relation \cite{hansen1982gap},
\begin{equation}
E_g(x,T)=-0.302+1.93x-0.81x^2+0.832x^3+5.35\times10^{-4}T(1-2x),
\end{equation}
with energy in electronvolts.

The above-bandgap absorption coefficient uses the Moazzami \emph{et al.} parameterization \cite{moazzami2005absorption},
\begin{align}
\alpha(E,x,T) &= K(x,T)\left(\frac{E-E_g}{E}\right)^{n(x,T)},\qquad E>E_g,\\
K(x,T) &= -20060+115750x+32.43T-64170x^2+0.43231T^2-101.92xT,\\
n(x,T) &= 0.74487-0.44513x+(0.000799-0.000757x)T,
\end{align}
and $\alpha=0$ for $E\le E_g$.  Here $\alpha$ is used in $\mathrm{cm^-1}$ in the optical-depth integral.  For a selected wavelength the unnormalized absorbed-generation density is
\begin{equation}
\tilde g(z)=\alpha(z)\exp\left[-\int_0^z\alpha(u)\,du\right],
\end{equation}
and the manuscript kernel is $g(z)=\tilde g(z)/\int_0^L\tilde g(u)du$.  Each wavelength is found numerically so that the normalized mean $\int zg(z)dz$ equals its target depth.

\begin{table}[h]
\caption{Theoretical generation channels supplied exactly to the forward average and inverse.  $P_{\rm abs}$ is the modeled absorbed probability through the $7.6\,\mu\mathrm{m}$ layer.}
\centering
\begin{tabular}{ccccc}
\toprule
Target mean ($\mu$m) & $\lambda$ ($\mu$m) & $P_{\rm abs}$ & Reintegrated mean ($\mu$m) & $\sigma_z$ ($\mu$m)\\
\midrule
2.000 & 2.059342009 & 0.999996080 & 2.000000000 & 0.783254\\
2.500 & 2.134650687 & 0.999984721 & 2.500000000 & 0.787048\\
3.000 & 2.215042347 & 0.999943407 & 3.000000000 & 0.789835\\
3.500 & 2.301173443 & 0.999801116 & 3.500000000 & 0.790930\\
4.000 & 2.393906801 & 0.999337640 & 4.000000000 & 0.788849\\
4.500 & 2.494502544 & 0.997909184 & 4.500000000 & 0.780666\\
\bottomrule
\end{tabular}
\end{table}

The optical model uses a dense one-dimensional grid of 12001 points.  The kernel construction contains no experimental optical-model uncertainty in the present stress; such uncertainty would be an additional nuisance in an experiment.

\section{Deterministic planar-depletion transport stress}

The central numerical counterexample is separate from the graded optical model above.  It solves a two-dimensional rectangular domain of lateral width $16\,\mu\mathrm{m}$ and absorber thickness $L=7.6\,\mu\mathrm{m}$.  The central case uses a full-width collecting contact and fixed endpoint potentials corresponding to $V_{\rm bias}=0.30\,\mathrm{V}$.  Over the collector-side region of width $W_d=3.0\,\mu\mathrm{m}$, the physical-potential solver imposes the Poisson curvature
\begin{equation}
\frac{d^2V}{dz^2}=\frac{2\Delta}{W_d^2},
\qquad \Delta=0.05\,\mathrm{V},
\end{equation}
while outside that region the planar continuum limit has $V''=0$.  The parameter $\Delta$ therefore controls curvature; it is not an independently added terminal voltage.  The characteristic curvature-field scale is $\Delta/W_d=166.7\,\mathrm{V/cm}$.

For the full-contact planar limit, let $a=L-W_d$.  The exact continuum solution consistent with $V(0)=0$ and $V(L)=V_{\rm bias}$ is
\begin{equation}
V(z)=
\begin{cases}
Az, & 0\le z<a,\\[3pt]
Az+\dfrac{\Delta}{W_d^2}(z-a)^2, & a\le z\le L,
\end{cases}
\qquad
A=\frac{V_{\rm bias}-\Delta}{L}.
\label{eq:planar-poisson}
\end{equation}
For the manuscript parameters, $A=0.0328947\,\mathrm{V/\mu m}$.  Accordingly, the electric-field magnitude across the curved region increases from $328.947$ to $662.281\,\mathrm{V/cm}$, with regional mean $495.614\,\mathrm{V/cm}$.  The corresponding uniform-bias field is $V_{\rm bias}/L=394.737\,\mathrm{V/cm}$, so the curved-region mean exceeds the uniform-bias value by $100.877\,\mathrm{V/cm}$.  The actual increase in potential drop across that $3.0\,\mu\mathrm{m}$ region relative to the uniform profile is $0.030263\,\mathrm{V}$, not $0.05\,\mathrm{V}$.  The physical electrostatic potential and the Shockley--Ramo weighting potential are solved independently.

Carrier transport is deterministic: $D_{\rm micro}=0$ and recombination is omitted.  The local drift speed is
\begin{equation}
v(E)=\frac{\mu E}{\sqrt{1+(\mu E/v_{\rm sat})^2}},
\end{equation}
with $\mu=0.90\,\mathrm{m^2/(V\,s)}$ and $v_{\rm sat}=6.0\times10^4\,\mathrm{m/s}$.  Along each deterministic trajectory the complex terminal-current transfer is accumulated as
\begin{equation}
H(\omega)=\int e^{-i\omega t}\,d\phi_w,
\end{equation}
including the final electrode contribution.  At dc this gives the exact Shockley--Ramo identity $H(0)=1-\phi_w(\mathbf r_0)$ for collected trajectories.

The manuscript baseline uses a $121\times91$ physical/weighting-potential mesh, $13$-point lateral Gaussian quadrature, $41$ source-depth nodes, and trajectory step $0.020\,\mu\mathrm{m}$.  The finite channel is then
\begin{equation}
J_m(\omega)=\int g_m(z)H(z,\omega)\,dz
\end{equation}
with the theoretical kernels above.  The homogeneous inverse fits the finite-kernel one-mode exponent and interprets it through $D\gamma^2+w\gamma=-i\omega$.

\section{Independent inferential convergence gate}

The convergence decision was declared before the successful run.  Mesh, source quadrature, and trajectory step were refined independently around the baseline.  The three levels were
\begin{align*}
\text{field mesh:}\quad &(81,61)\rightarrow(121,91)\rightarrow(161,121),\\
\text{source quadrature:}\quad &(9,31)\rightarrow(13,41)\rightarrow(17,61),\\
\text{trajectory step:}\quad &0.035\rightarrow0.020\rightarrow0.0125\ \mu\mathrm{m}.
\end{align*}
Seven unique configurations were required because all axes share the same baseline.

The predeclared baseline-to-fine tolerances were $2\%$ relative for probe-frequency $D_{\rm eff}$, $0.5\%$ for $w_{\rm eff}$, $2\%$ and $0.5\%$ for the low-band joint $D,w$ fit, $0.002$ absolute for the 1-GHz homogeneous-law residual, $2\times10^{-5}$ absolute for the maximum one-mode kernel-fit residual through 1 GHz, $1\%$ of the fine finite-kernel 100-MHz diffusion scale for the upstream point-control change, $2\%$ relative for the inside-region point-control diffusion, $10^{-10}$ for the dc Ramo error, and a minimum collection fraction of 0.999.  Positive-$D$ sign stability was also required for the finite-kernel 100/500/1000-MHz inversions and the inside-region point-source control.

\begin{table}[h]
\caption{Field-mesh convergence of the same-frequency finite-kernel inverse.  The field mesh is the limiting numerical coordinate.}
\centering
\begin{tabular}{ccccc}
\toprule
Frequency & Coarse $D_{\rm eff}$ & Baseline $D_{\rm eff}$ & Fine $D_{\rm eff}$ & Baseline$\to$fine\\
 & \multicolumn{3}{c}{($\mathrm{m^2/s}$)} & (\%)\\
\midrule
100 MHz & 0.002526515 & 0.002609795 & 0.002653537 & 1.648\\
500 MHz & 0.002467684 & 0.002548603 & 0.002591124 & 1.641\\
1 GHz   & 0.002274809 & 0.002348945 & 0.002387949 & 1.633\\
\bottomrule
\end{tabular}
\end{table}

The corresponding baseline-to-fine drift-speed changes are 0.180\%, 0.176\%, and 0.163\%.  The low-band joint $D$ and $w$ changes are 1.645\% and 0.180\%.  The 1-GHz law residual changes by 1.752e-04 absolute and the maximum one-mode finite-kernel fit residual through 1 GHz changes by 3.361e-06 absolute.

Source/kernel quadrature is substantially more converged: the baseline-to-fine relative changes in $D_{\rm eff}$ are 3.332e-05, 3.426e-05, and 3.758e-05 at 100, 500, and 1000 MHz.  Trajectory-step changes are 4.995e-06, 5.738e-06, and 8.208e-06.

The causal point controls remain separated.  The upstream point-source sequence gives $D_{\rm eff}=1.735e-12\,\mathrm{m^2/s}$ at baseline and $9.174e-12\,\mathrm{m^2/s}$ on the fine field mesh, both numerical-zero scale compared with the finite-kernel result.  The point-source sequence entirely inside the nonuniform region gives $4.871054e-03\,\mathrm{m^2/s}$ at baseline and $4.868700e-03\,\mathrm{m^2/s}$ on the fine mesh.  All seven numerical configurations satisfy the collection and dc Shockley--Ramo integrity checks, and every predeclared convergence gate passes.

A convergence PASS establishes numerical stability of this deterministic surrogate and inverse under the declared refinements.  It does not establish experimental feasibility, device calibration, or uniqueness of the physical interpretation.

\section{Reproduction record}

The principal executable files for this supplement are:
\begin{itemize}
\item \texttt{numerics/hgcdte\_ramo\_four\_color\_gradient\_prediction.py} --- theoretical HgCdTe optical-kernel generator;
\item \texttt{numerics/realistic\_geometry\_closure\_stress.py} --- two-dimensional physical and weighting potentials plus deterministic Shockley--Ramo trajectories;
\item \texttt{numerics/paper02\_kernel\_aware\_depletion\_frequency\_law.py} --- finite-kernel homogeneous transport inverse;
\item \texttt{numerics/paper02\_point\_vs\_kernel\_causal\_test.py} --- upstream and inside-region point-source controls;
\item \texttt{numerics/paper02\_inference\_convergence\_gate.py} and \texttt{paper02\_inference\_convergence\_runner.py} --- independent numerical-refinement gate.
\end{itemize}

The corrected checked convergence execution is GitHub Actions run \texttt{31948607702}, job \texttt{95168474631}, artifact \texttt{paper02-inference-convergence} (artifact id \texttt{9264012168}).  The downloaded artifact SHA-256 recorded at validation was \texttt{ac9c9ade5e658fedd6ff846ee869dcc25d5bb0e4d85dd26a9657e6cf3dfaf275}.  The immediately preceding execution completed all seven numerical solves but failed in the reporting layer because the low-band metric label was mis-dispatched as a frequency label; it produced no convergence verdict.  That failed route is retained as provenance.

\bibliographystyle{apsrev4-2}
\bibliography{PAPER02_REFERENCES_REV4}

\end{document}
